%%
%% Copyright 2007-2024 Elsevier Ltd
%%
%% This file is part of the 'Elsarticle Bundle'.
%% ---------------------------------------------
%%
%% It may be distributed under the conditions of the LaTeX Project Public
%% License, either version 1.3 of this license or (at your option) any
%% later version.  The latest version of this license is in
%%    http://www.latex-project.org/lppl.txt
%% and version 1.3 or later is part of all distributions of LaTeX
%% version 1999/12/01 or later.
%%
%% The list of all files belonging to the 'Elsarticle Bundle' is
%% given in the file `manifest.txt'.
%%
%% Template article for Elsevier's document class `elsarticle'
%% with numbered style bibliographic references
%% SP 2008/03/01
%% $Id: elsarticle-template-num.tex 249 2024-04-06 10:51:24Z rishi $
%%
\documentclass[preprint,12pt]{elsarticle}

%% Use the option review to obtain double line spacing
%% \documentclass[authoryear,preprint,review,12pt]{elsarticle}

%% Use the options 1p,twocolumn; 3p; 3p,twocolumn; 5p; or 5p,twocolumn
%% for a journal layout:
%% \documentclass[final,1p,times]{elsarticle}
%% \documentclass[final,1p,times,twocolumn]{elsarticle}
%% \documentclass[final,3p,times]{elsarticle}
%% \documentclass[final,3p,times,twocolumn]{elsarticle}
%% \documentclass[final,5p,times]{elsarticle}
%% \documentclass[final,5p,times,twocolumn]{elsarticle}

%% For including figures, graphicx.sty has been loaded in
%% elsarticle.cls. If you prefer to use the old commands
%% please give \usepackage{epsfig}
%% The amssymb package provides various useful mathematical symbols
\usepackage{amssymb}
%% The amsmath package provides various useful equation environments.
\usepackage{amsmath,amsthm,proof,float}
\usepackage{semantic}
\usepackage{mathtools}
\usepackage{listings}
\usepackage{hyperref}
\def\lstlanguagefiles{lstlean.tex}
% set default language
\lstset{language=lean}
\newtheorem{Lemma}{Lemma}
\newtheorem{Theorem}{Theorem}
\theoremstyle{definition}
\newtheorem{Example}{Example}

%% The amsthm package provides extended theorem environments
%% \usepackage{amsthm}

%\usepackage{xcolor}

\usepackage{color}
\definecolor{keywordcolor}{rgb}{0.7, 0.1, 0.1}   % red
\definecolor{tacticcolor}{rgb}{0.0, 0.1, 0.6}    % blue
\definecolor{commentcolor}{rgb}{0.4, 0.4, 0.4}   % grey
\definecolor{symbolcolor}{rgb}{0.0, 0.1, 0.6}    % blue
\definecolor{sortcolor}{rgb}{0.1, 0.5, 0.1}      % green
\definecolor{attributecolor}{rgb}{0.7, 0.1, 0.1} % red

\def\lstlanguagefiles{lstlean.tex}
% set default language
\lstset{language=lean}

\usepackage{lineno}
\usepackage{url}

%% The lineno packages adds line numbers. Start line numbering with
%% \begin{linenumbers}, end it with \end{linenumbers}. Or switch it on
%% for the whole article with \linenumbers.
%% \usepackage{lineno}

\journal{Journal of Computer Languages}

\begin{document}

\begin{frontmatter}

%% Title, authors and addresses

%% use the tnoteref command within \title for footnotes;
%% use the tnotetext command for theassociated footnote;
%% use the fnref command within \author or \affiliation for footnotes;
%% use the fntext command for theassociated footnote;
%% use the corref command within \author for corresponding author footnotes;
%% use the cortext command for theassociated footnote;
%% use the ead command for the email address,
%% and the form \ead[url] for the home page:
%% \title{Title\tnoteref{label1}}
%% \tnotetext[label1]{}
%% \author{Name\corref{cor1}\fnref{label2}}
%% \ead{email address}
%% \ead[url]{home page}
%% \fntext[label2]{}
%% \cortext[cor1]{}
%% \affiliation{organization={},
%%             addressline={},
%%             city={},
%%             postcode={},
%%             state={},
%%             country={}}
%% \fntext[label3]{}

\title{Pest control: A formal model of the Pest parser generator}
%% use optional labels to link authors explicitly to addresses:
%% \author[label1,label2]{}
%% \affiliation[label1]{organization={},
%%             addressline={},
%%             city={},
%%             postcode={},
%%             state={},
%%             country={}}
%%
%% \affiliation[label2]{organization={},
%%             addressline={},
%%             city={},
%%             postcode={},
%%             state={},
%%             country={}}

\author[ufop]{Guilherme Daher} %% Author name
\author[ufop]{Elton Cardoso}
\author[ufsc]{Samuel Feitosa}
\author[ufop]{Rodrigo Ribeiro} %% Author name

%% Author affiliation
\affiliation[ufop]{organization={Federal University of Ouro Preto},%Department and Organization
            city={Ouro Preto},
            state={MG},
            country={Brazil}}
\affiliation[ufsc]{organization={Federal University of the Southern Border},%Department and Organization
            city={Chapecó},
            state={SC},
            country={Brazil}}

%% Abstract
\begin{abstract}
Parsing expressions grammars (PEGs) are a recognition-based formalism for parsers,
which has been gaining popularity because they avoid some problems present in context-free grammar
based specifications. In the context of the Rust programming language, the most used
PEG-based parser generator is Pest. An interesting feature of Pest grammars is that they
support new operators for repetition and handling an implicit stack during parsing.
In this work, we propose a formalization of these new constructs and extend a
type inference algorithm which guarantee that all well-typed Pest-style grammars
terminate on all inputs.
\end{abstract}

%%Graphical abstract
%\begin{graphicalabstract}
%\includegraphics{grabs}
%\end{graphicalabstract}

%%Research highlights
\begin{highlights}
\item Formalizes the semantics of Pest parser generator new repetition and stack operators
  using an operational semantics.
\item Extends an existing PEG type system and its inference algorithm that ensures
  termination for all well-typed grammars to consider Pest new operators.
\item We provide a Lean 4 mechanized proof of type soundness for the defined type system which
  ensures termination of parsing for all well-typed grammars.
\end{highlights}

%% Keywords
\begin{keyword}
Parsing \sep parsing expression grammars \sep type systems
\end{keyword}

\end{frontmatter}

%% Add \usepackage{lineno} before \begin{document} and uncomment
%% following line to enable line numbers
\linenumbers

%% main text
%%

\section{Introduction}\label{sec:intro}

Parsing is a core problem in computer science. In simple terms, a parsing
algorithm, also named a parser, checks if a string of symbols conforms to some
set of rules specified by formalisms like context-free grammars (CFGs) or
regular expressions (REs). Usually, parsers construct data structures showing
what rules were used to conclude that the input string was obtainable from the
specification or an error message indicating that such string is unobtainable
from the specification. Due to its importance, parsing is an extensively studied
subject, as evidenced by Grune and Jacob's book "Parsing Techniques: A Practical
Guide", which has more than 600 pages and covers only classic algorithms \cite{Grune90}.
However, parsing is far from being a solved problem: in recent years,
we have seen a growing interest in the development of new algorithms and
formalisms~\cite{Ford2004, Trevor10,Asperti10, Lucks17, Zhang23} or the verification
of well-known parsing techniques ~\cite{FirsovU14, FirsovU13,LopesRC16,Bernardy16}.

In the early 2000's, Brian Ford proposed \emph{Parsing Expression Grammars}
(PEGs),  a recognition-based formalism to
define languages~\cite{Ford2004}. The novelty of PEGs is that they
specify how to check whether a string is in
the language being defined and not how it is generated by its specification.
Formally a PEG is a 4-tuple $(V_N, V_T, R, e_S)$ where $V_N$ is the set of non-terminals
(a.k.a. variables), $V_T$ is the set of terminals (a.k.a. alphabet symbols),
$R$ is a function which maps a non terminal from $V_N$ to a {\it parsing expression}
and $e_S$ is the initial parsing
expression. A parsing expression is very
similar to the right-hand side of a CFG in the extended Backus-Naur
form with the addition of new operators such as the not-predicate (!) and
the prioritized choice (/). A fundamental difference between PEG
and CFGs is the choice operator, which impose priority between alternatives.
It means that $\beta$ in a prioritized choice expression $\alpha / \beta$ is tested only
in case of failure on $\alpha$. The not-predicate operator checks
if the string matches some syntax without consuming the input.

In the context of the Rust programming language~\cite{RustBook}, the
Pest library~\cite{PEST} is the most popular PEG-based parser
generator for Rust, having more than 114.000 downloads registered at
\href{http://crates.io}{crates.io} package registry. Besides supporting
all traditional PEG operators, Pest support operators to manipulate an
implicit stack during parser execution. The inclusion of a stack allows
Pest parsers to elegantly express identation sensitive languages by
controlling identation levels using the stack. While such new operators
eases the task of writing grammars for complex formats, additional care
must be taken in order to ensure PEG \emph{well-formedness}. We
say that a PEG is well-formed if it terminates any input
strings~\cite{Ford2004}. Pest users have reported that the tool
fails to detect grammars that loop on some inputs~\cite{PESTBUG}.
Aiming to solve this problem, we proposed a formalization of a generalized
version of these Pest operators using operational semantics and
an extended version of a type system for PEGs~\cite{Ribeiro19} and its
inference algorithm~\cite{Cardoso23} to deal with these operators.

Specifically, we make the following contributions:

\begin{itemize}
  \item We formalize both stack and the new repetitions Pest operators using
    operational semantics and an extended version of a type system for PEGs~\cite{Ribeiro19}
    to deal with these new operations. The proposed semantics is realized by an
    intrisincally-typed definitional interpreter~\cite{Poulsen17} and we prove
    a type soundness theorem~\cite{Amin17} for it using the Lean 4 proof assistant~\cite{demoura2021}.
    The complete formalization source code is available on-line~\cite{pegrepository}.
  \item We include operations to query the top-most element of the stack
    which can be used to handle context-sensitive formats, like length
    indexed fields, common in binary data formats;
  \item We show how to extend a type inference algorithm~\cite{Cardoso23} for PEGs
    to deal with Pest's new operators. A prototype implementation of the inference
    algorithm was developed using the Racket programming language~\cite{Culpepper19}
    and used to check some example grammars.
\end{itemize}

This paper extends our previous work~\cite{Daher25} by: 1) giving details of the
formalization of the type system and the proposed semantics in Lean 4; 2)
a certified proof of a type soundness theorem for Pest-style grammars and 3)
we give additional examples of how our generalized semantics can be used to
express complex formats in an elegant way.

The rest of this paper is organized as follows.
Section~\ref{sec:pest-operators} presents an informal description of Pest operators.
Section~\ref{sec:background}
reviews the syntax, semantics, a type system for
PEGs and provides an overview of the some features of Lean 4 used in our formalization.
Section~\ref{sec:pest-formalization} defines the syntax, semantics, the type
system and its inference algorithm for the new PEG operators. Section~\ref{sec:typesoundnesslean}
describes our formalization of the type soundness theorem for Pest grammars using Lean 4.
Section~\ref{sec:examples}
presents some examples of grammars accepted by our tool.
Related work is discussed on Section~\ref{sec:related} and
Section~\ref{sec:conclusion} concludes.


\section{An overview of Pest new operators}~\label{sec:pest-operators}

Before starting the description of each new operator, we describe
\emph{repetition indexes}, which are values ranged over natural numbers and a
special value, \textbf{infty}, which denotes an upper-bound over
indexes\footnote{For any index $i \neq \mathbf{infty}$, we have
that $i < \mathbf{infty}$ and $i \circ \mathbf{infty} = \infty$}.
Notation $\circ$ denotes either addition
or multiplication operations over natural numbers. The constants
\textbf{top.tonat} and \textbf{top.length} refer to the conversion
of the stack top element to a numeric value and its size, respectively.
More details on how these operations work will be described
on the semantics.

The first operator we consider is the iterated repetition, denoted $e^i$ in
the abstract syntax and \verb|e^i| in the concrete syntax. It describes that
expression $e$ should be repeated $i$ times over the input, failing if any of these iterations
results in a parser failure. As an example of the usefulness of this operator, consider
the task of representing a PEG for a format formed by two bytes values followed by
a single bit terminator. We could express this simple format as the following
grammar using the iterated repetition:
\begin{verbatim}
bit <-- '0' / '1'
byte <-- bit^8
start: byte^2 ~ bit
\end{verbatim}
Notation \verb|'0'| denotes an expression for a string,
operator \verb|~| denotes concatenation (following Pest concrete syntax),
\verb|A <-- e| denotes a grammar rule and \verb|start:| specifies
the starting expression for this grammar. The usage of the iterated
repetition operator allows for concise description of pattern ``2 bytes''
as simply \verb|byte^2|, which indicates that non terminal \verb|byte|
should be repeated two times over the input.

The Pest tool also introduces an interval repetition operator, $e[i,j]$,
which executes the expression $e$ a number of times between the indexes $i$
and $j$, inclusive. As an example of this operator usage, consider the
task of parsing either 2, 3 or 4 bits sequences from input. This could be
expressed elegantly by the following expression using the interval repetition:
\verb|('0' / '1')[2,4]|.

However, not all combinations of indexes are allowed in repetition operators.
An example of an invalid repetition operator: \verb|e[infty, 2]|, since a
bounded repetition, \verb|e[i,j]|, is valid if $i < j$ and $i$ is not the constant
\verb|infty|. Regarding the exact repetition operator, $e^i$, index $i$ is not allowed to
be \textbf{infty}, since its semantics is to parse the input by repeating
an expression $e$ \textbf{exactly} $i$ times, which makes the expression
$e^{\mathbf{infty}}$ to be repeat indefinitely over the input.

Now, we turn our attention to the stack manipulation operators. Expression
\verb|push(e)| parses the input using the expression \texttt{e} and pushes the
consumed input on the parser stack. Operator \verb|pop| removes the current
stack top-most element and tries to match it on the current input,
\verb|peek| also matches the input without removing the element from the
stack and \verb|drop| removes the first element of current stack, without
matching it. The operator \verb|peekall| tries to match the input using the
complete stack content, without emptying the stack. Finally, \verb|dropall|
empties the stack.

One interesting example use of stack operators is to parse
length-indexed fields present in several data formats.
The next grammar shows how to express netstrings, which is a format method
for strings which contains a field to denote the size of the represented data.
\begin{verbatim}
number <-- '0' / '1' / ... / '9'
letter <-- 'a' / ... / 'z'
char <-- number / letter

start: push(number+) ~ ':' ~ char^top.tonat ~ ','
\end{verbatim}
The starting expression begins by pushing the input size value on the parser stack,
and then access this size information using \verb|top.tonat| to parse the exact
quantity of input based on the size information that starts the netstring.
Operators \verb|top.tonat| and \verb|top.length| are not present in the current
definition of Pest parser generator. We include both since they have a direct
formalization and increases the expressivity for Pest grammars which can be used
to represent length-indexed fields, which are notoriously difficult to parse
correctly~\cite{Zhang23, Lucks17}.


\section{Background}\label{sec:background}

\paragraph{An Overview of PEGs}

Intuitively, PEGs are a formalism for describing top-down parsers.
Formally, a PEG is a 4-tuple $(V,\Sigma,R,e_S)$, where $V$ is a
finite set of variables, $\Sigma$ is the alphabet, $R$ is the finite
set of rules, and $e_s$ is the start expression. Each rule $r \in R$ is a
pair $(A,e)$, usually written $A \leftarrow e$, where $A \in V$ and $e$
is a parsing expression. We let the meta-variable $a$ denote an
arbitrary alphabet symbol, $A$ a variable and $e$ a parsing expression.
Following common practice, all meta-variables can appear primed or
sub-scripted. The following context-free grammar defines
the syntax of a parsing expression:
\[
\begin{array}{lcl}
e & \to  & \epsilon \,\mid\, a \,\mid\, A\,\mid\,e_1\:e_2\,
\mid\,e_1\,/\,e_2\,\mid\,e^\star\,\mid\,!\,e\\
\end{array}
\]
The execution of parsing expressions is defined by an inductively defined
judgment that relates pairs formed by a parsing expression and an input string
to pairs formed by the consumed prefix and the remaining string.
Notation $(e,s) \Rightarrow_G (s_p,s_r)$ denote that parsing expression $e$
consumes the prefix $s_p$ from the input string $s$ leaving the suffix $s_r$.
The notation $(e,s) \Rightarrow_G \bot$ denote the fact that $s$ cannot be parsed by
$e$. We let meta-variable $r$ denote an arbitrary parsing result, i.e.,
either $r$ is a pair $(s_p,s_r)$ or $\bot$. We say that an expression $e$
fails if its execution over an input produces $\bot$; otherwise, it succeeds.
Figure~\ref{fig:pegsemantics} defines the PEG semantics.
\begin{figure*}[h!]
   \[
      \begin{array}{cc}
         \infer[_{\{Eps\}}]{(\epsilon,s)\Rightarrow_G (\epsilon,s)}{} &
         \infer[_{\{ChrS\}}]{(a,as_r) \Rightarrow_G (a,s_r)}{} \\ \\
         \infer[_{\{ChrF\}}]{(a,bs_r) \Rightarrow_G \bot}{a \neq b} &
         \infer[_{\{Var\}}]{(A,s) \Rightarrow_G r}
                           {A \leftarrow e \in R & (e,s) \Rightarrow_G r} \\ \\
         \infer[_{\{Not_S\}}]{(!\,e,s) \Rightarrow_G (\epsilon,s)}
         {(e,s) \Rightarrow_G \bot}
           &
         \infer[_{\{ChrNil\}}]{(a,\epsilon) \Rightarrow_G \bot}{}\\ \\

         \multicolumn{2}{c}{
            \infer[_{\{Cat_{S1}\}}]{(e_1\,e_2,s_{p_1}s_{p_2}s_r) \Rightarrow_G (s_{p_1}s_{p_2}, s_r)}
                                 {(e_1,s_{p_1}s_{p_2}s_r) \Rightarrow_G (s_{p_1},s_{p_2}s_r) &
                                 (e_2,s_{p_2}s_r)\Rightarrow_G (s_{p_2},s_r)}
         } \\ \\
         \multicolumn{2}{c}{
            \infer[_{\{Cat_{F2}\}}]{(e_1\,e_2,s_ps_r) \Rightarrow_G \bot}
                                 { (e_1,s_ps_r) \Rightarrow_G (s_p,s_r) &
                                    (e_2,s_r) \Rightarrow_G \bot}} \\ \\
         \infer[_{\{Cat_{F1}\}}]{(e_1\,e_2,s)\Rightarrow_G \bot}{(e_1,s) \Rightarrow_G \bot} &
         \infer[_{\{Alt_{S1}\}}]{(e_1\,/\,e_2,s_p\,s_r) \Rightarrow_G (s_p,s_r)}
                                {(e_1,s_p\,s_r)\Rightarrow_G (s_p,s_r)}\\ \\
         \multicolumn{2}{c}{
            \infer[_{\{Alt_{S2}\}}]{(e_1\,/\,e_2,s_p\,s_r) \Rightarrow_G r}
                                  {(e_1,s_p\,s_r)\Rightarrow_G \bot &
                                   (e_2,s_p\,s_r)\Rightarrow_G r}
         } \\ \\
         \multicolumn{2}{c}{
            \infer[_{\{Star_{rec}\}}]{(e^\star,s_{p_1}s_{p_2}s_r) \Rightarrow_G (s_{p_1}s_{p_2},s_r)}
                                 {(e,s_{p_1}s_{p_2}s_r) \Rightarrow_G (s_{p_1},s_{p_2}s_r) &
                                  (e^\star, s_{p_2}s_r) \Rightarrow_G (s_{p_2},s_r)}
         }\\ \\
         \multicolumn{2}{c}{
            \infer[_{\{Star_{end}\}}]{(e^\star,s) \Rightarrow_G (\epsilon,s)}
                                    {(e,s) \Rightarrow_G \bot}} \\ \\
         \multicolumn{2}{c}{
            \infer[_{\{Not_F\}}]{(!\,e,s_p\,s_r) \Rightarrow_G \bot}
                               {(e,s_p\,s_r) \Rightarrow_G (s_p,s_r)}
         }
     \end{array}
   \]
   \centering
   \caption{Parsing expressions operational semantics.}
   \label{fig:pegsemantics}
\end{figure*}
We comment on some rules of the semantics. Rule $_{Eps}$ specifies that
expression $\epsilon$ will not fail on any input $s$ by leaving it unchanged.
Rule $_{ChrS}$ specifies that an expression $a$ consumes the first character when
the input string starts with an `a' and rule $_{ChrF}$ shows that it fails when
the input starts with a different symbol. Rule $_{Var}$ parses the input using
the expression associated with the variable in the grammar $G$. When parsing a
sequence expression, $e_1\:e_2$, the result is formed by $e_1$ and $e_2$ parsed
prefixes and the remaining input is given by the result $e_2$. Rules $_{Cat_{F1}}$ and
$_{Cat_{F2}}$ say that if $e_1$ or $e_2$ fail, then the whole expression fails.
The rules for choice impose that we only try
expression $e_2$ in $e_1 / e_2$ when $e_1$ fails. Parsing a star
expression $e^\star$ consists in repeatedly execute $e$ on the input string.
When $e$ fails, $e^\star$ succeeds without consuming any symbol of the input
string. Finally, the rules for the not-predicate expression, $!\,e$, specify
that whenever the expression $e$ succeeds on input $s$, $!\,e$ fails; and when $e$
fails on $s$ we have that $!\,e$ succeeds without consuming any input.

\paragraph{A type system for PEGs}

A key issue in the PEG formalism is to ensure the termination of the parsing algorithm
for arbitrary inputs. Ford~\cite{Ford2004} proposed a well-formedness predicate
which is computed as fixed-point over the set of grammar's sub-expressions.
In simple terms, a PEG can loop indefinitely over an input string if: 1) it has
direct or indirect left-recursive rules ; 2) expressions that succeed without
consuming any symbol under the Kleene star operator.
As pointed by~\cite{Ribeiro19}, one inconvenience of Ford's well-formedness predicate
is that it is not compositional.
Based on this observation, the work~\cite{Ribeiro19} propose a type
system that guarantees termination by restricting valid expressions to those
that do not obey conditions 1 and 2. In order to ensure these restrictions,
types are interpreted as records containing two fields:
a boolean value indicating that the current expression may succeed without
consuming anything from the input string and a set of variables that can
appear as the first symbol on the recognition branch from the current parsing
expression. This set (named head-set) is defined by recursion on the parsing
expression syntax as follows:
\[
  \begin{array}{lcl}
    head (\epsilon) & = & \emptyset \\
    head (a) & = & \emptyset \\
    head (A) & = & \{A\} \cup head(R(A)) \\
    head (e_1\,e_2) & = & \left\{
                          \begin{array}{ll}
                            head(e_1) \cup head(e_2) & \text{if }e_1 \rightharpoonup 0\\
                            head(e_1) & \text{otherwise}
                          \end{array}
                                        \right. \\
    head(e_1\,/\,e_2) & = & head(e_1) \cup head(e_2) \\
    head(e^\star) & = & head(e) \\
    head(!\,e)    & = & head(e)
  \end{array}
\]
Basically, the head-set of $\epsilon$ and $a \in \Sigma$ is the empty set.
For a non terminal, the head-set is formed by the non terminal itself plus
the head-set of its associated parsing expression. The head-set of a sequence
expression $e_1\:e_2$ need to include the head-set of $e_2$, whenever $e_1$
succeeds without consuming anything ($e_1 \rightharpoonup 0$). Otherwise,
$head(e_1\:e_2)$ is just $head(e_1)$. The set of non terminals for the choice
operator is the union of the head-set of its operands. Finally, the head-set
for the Kleene star and not operator are just the set for its underlying expressions.
Following~\cite{Ribeiro19}, $\tau$ denotes an arbitrary type, $\tau.null$
denotes the boolean field of $\tau$, and $\tau.head$ the set of
variables that can appear as the first symbol of $\tau$'s parsing expression.
We let notation $\langle b , S \rangle$ denote a type $\tau$ formed by a
boolean $b$ and a set $S$. In the type system~\cite{Ribeiro19}, we use the following operations:
\[
  \begin{array}{l}
    b \Rightarrow S = \text{if } b \text{ then } S \text{ else }\emptyset\\
    \tau_1 \otimes \tau_2 = \langle \tau_1.null \land \tau_2.null,
                                \tau_1.head \cup (\tau_1.null \Rightarrow \tau_2.head) \rangle \\
    \tau_1 \oplus \tau_2 = \langle \tau_1.null \lor \tau_2.null,
                                \tau_1.head \cup \tau_2.head\rangle\\
    !\,\tau = \langle true , \tau.head \rangle\\
    \langle false, S \rangle^\star = \langle true, S \rangle
  \end{array}
\]
The first operation is a boolean test that returns the set $S$
whenever the condition is true. The product operation ($\tau_1 \otimes \tau_2$) combines two
types by taking the conjunction of their boolean flags and the union of their respective
head-sets in case $\tau_1.null$ is true. The co-product operation ($\tau_1 \oplus \tau_2$) on types is similar to the
product, but instead of using conjunction it takes the disjunction of their boolean fields.
The not operation on types just change its boolean field to true and keep its head-set.
The last operation is the Kleene star on types, which also sets its boolean field to true.
However, the Kleene is only defined for types $\tau$ such that $\tau.null = false$, being
undefined when $\tau.null = true$. Finally, following the common practice, the meta-variable
$\Gamma$ denotes a typing context and notation $\Gamma(A) = \tau$ holds whenever
$(A,\tau) \in \Gamma$.
The type system is defined as an inductive judgment
$\Gamma \vdash e : \tau$ in Figure~\ref{fig:typesystem}. A
PEG $G = (V,\Sigma,R,e_S)$ is considered well-typed, written $\Gamma \vdash G$,
if $\forall A \in V . R(A) : \Gamma (A)$.

\begin{figure}[h!]
\[
  \begin{array}{cc}
    \infer{\Gamma \vdash \epsilon : \langle true , \emptyset \rangle}{} &
    \infer{\Gamma \vdash a : \langle false , \emptyset\rangle}{} \\ \\
    \multicolumn{2}{c}{
      \infer{\Gamma \vdash A : \langle \tau.null, \tau.head \cup {A}\rangle}
            {\Gamma(A) = \tau & A \not\in \tau.head}
    }\\ \\
    \infer{\Gamma\vdash e_1\,/\,e_2 : \tau_1 \oplus \tau_2}
          {\Gamma \vdash e_1 : \tau_1 & \Gamma \vdash e_2 : \tau_2}
    &
    \infer{\Gamma\vdash !\,e :\, !\,\tau}
          {\Gamma \vdash e : \tau}\\ \\
    \infer{\Gamma\vdash e_1\,e_2 : \tau_1 \otimes \tau_2}
          {\Gamma \vdash e_1 : \tau_1 & \Gamma \vdash e_2 : \tau_2} &
    \infer{\Gamma\vdash e^\star : \tau^\star}
          {\begin{array}{c}
             \Gamma \vdash e : \tau\\
             \tau.null = false\\
           \end{array}}
  \end{array}
\]
\centering
\caption{Type system definition.}
\label{fig:typesystem}
\end{figure}

The first rule of the type system specifies that a parsing
expression $\epsilon$ has type $\langle true, \emptyset \rangle$.
The rule of a single character expression says that its type is
$\langle false, \emptyset\rangle$, since it does not succeed without
consuming a symbol and no variable can appear as the head of its
recognition branch, i.e., its head-set is $\emptyset$. The type of a prioritized
choice expression ($e_1/e_2$) has type $\tau_1 \oplus \tau_2$,
where $\Gamma \vdash e_1 : \tau_1$ and $\Gamma \vdash e_2 : \tau_2$.
The type of a sequence expression is given by the product operation
on their sub-expressions types, and the star operation is only well-typed
if its underlying expression does not succeed without consuming
a symbol. A not-predicate parsing expression always has the type
$\langle true,\tau.head \rangle$ whenever its sub-expression is well-typed.
Finally, a variable is well-typed only if its type is in the context and
itself does not belong to its head-set.
The type system proposed by~\cite{Ribeiro19} readily rejects invalid expressions like
$a(!b / c)^\star$, since it only considers well-typed a Kleene parsing
expressions in which an underlying expression does
not succeed without consuming the input. However, an inconvenience of it is
the need to require type annotations on all grammar's non terminals.
A later work~\cite{Cardoso23} proposes a type inference algorithm, which
guarantees termination of the parsing process.

\paragraph{An overview of the Lean 4 proof assistant}
The Lean proof assistant started as a Microsoft Research
open source project in 2013. Lean is based on the
calculus of inductive constructions~\cite{Nederpelt14} and it supports
inductive types and functions, type classes, macros and an efficient
C backend~\cite{demoura2021}. The Lean proof assistant has been used
with success to formalize several mathematical theories in the
Mathlib project~\cite{Mathlib}.

In this work, we use the Lean 4 proof assistant to prove a
type soundness result for the type system developed for
Pest-style grammars. In order to present the main features
of Lean used in the formalization, we will present the
encoding of a type soundness result for a simple language of
arithmetic expressions, proposed in~\cite{Pierce02}. We
start by defining its syntax using inductive types, as follows:

\begin{lstlisting}
inductive Ty : Type where
| NAT | BOOL

inductive Exp : Ty → Type where
| True : Exp .BOOL
| False : Exp .BOOL
| Zero : Exp .NAT
| Succ : Exp .NAT → Exp .NAT
| Pred : Exp .NAT → Exp .NAT
| IsZero : Exp .NAT → Exp .BOOL
| If : ∀ {t}, Exp .BOOL → Exp t → Exp t → Exp t
\end{lstlisting}

An inductive type defines a new type name, its sort,
and a set of its data constructors. In the previous code
piece, we have two inductive types: one for representing
types, \lstinline|Ty|, and another for expressions, \lstinline|Exp|.
The sort \lstinline|Type| is the "type" of usual datatypes in Lean
and forms a hierarchy indexed by levels, where types of a lower level
are contained in higher ones. In the previous definition, inductive type
\lstinline|Ty| is defined in sort \lstinline|Type| and consists of two
constructors representing the base types of the expression language:
\lstinline|NAT| for natural numbers and \lstinline|BOOL| for booleans.
The \lstinline|Exp| type follows the so-called intrinsically-typed
representation of programs~\cite{Poulsen17,Feitosa22}, which ensures that only
well-typed programs can be constructed. Specifically, \lstinline|Exp| is
indexed by \lstinline|Ty|, meaning that each expression constructor is
associated with its type. The boolean constructors \lstinline|True| and
\lstinline|False| produce expressions of type \lstinline|Exp .BOOL|,
while \lstinline|Zero| produces an expression of type \lstinline|Exp .NAT|.
The \lstinline|Succ| and \lstinline|Pred| constructors are functions that
take a natural number expression and return a natural number expression,
enforcing that these operations apply only to naturals. The
\lstinline|IsZero| constructor maps from \lstinline|Exp .NAT| to
\lstinline|Exp .BOOL|, representing an equal-to-zero test for numeric values.
Finally, the \lstinline|If| constructor takes a boolean condition and ensures that
its two branches have the same type, \lstinline|Exp t|. This intrinsic typing
discipline renders ill-typed terms inexpressible, guaranteeing type safety by
construction.

Following~\cite{Amin17}, the next ingredient in our type soundness proof is
the definition of an interpreter for the expression language. Before presenting
the interpreter, we need to define the semantic values which will be the
result of evaluating expressions. Values are either boolean constants,
zero or the successor of a numeric value. We denote values by the following
inductive types.
\begin{lstlisting}
inductive Val : Ty → Type where
| True : Val .BOOL
| False : Val .BOOL
| Zero : Val .NAT
| Succ : Val .NAT → Val .NAT
\end{lstlisting}
The intepreter is defined by the \lstinline|interp| function which is a
fuel-based interpreter for expressions, mapping each well-typed
expression \lstinline|e : Exp t| to an optional value
\lstinline|Option (Val t)|. The only situation in which the intepreter
produces a \lstinline|Option.none| result is when the fuel counter reachs
zero. The usage of fuel based interpreters is a standard technique for
implementing semantics in proof assistants~\cite{Amin17}.

\begin{lstlisting}
def interp {t}(fuel : ℕ)(e : Exp t) : Option (Val t) :=
  match fuel with
  | 0 => .none
  | fuel1 + 1 =>
    match e with
    | .True => .some Val.True
    | .False => .some .False
    | .Zero => .some .Zero
    | .Succ e1 => do
        let v ← interp fuel1 e1
        vsucc v
    | .Pred e1 => do
        let v ← interp fuel1 e1
        vpred v
    | .IsZero e1 => do
        let v ← interp fuel1 e1
        viszero  v
    | .If e1 e2 e3 => do
        let v ← interp fuel1 e1
        match v with
        | .True => interp fuel1 e2
        | .False => interp fuel1 e3
\end{lstlisting}

The interpreter starts by pattern matching on the fuel argument, returning
\lstinline|.none|\footnote{It possible to omit the type name in data
constructor expressions, when Lean is able to infer its correct type
from the expression context.}, when \lstinline|fuel = 0|.
For non-zero fuel values, the function pattern matches on the expression
structure. The base cases for boolean and natural number literals
(\lstinline|True|, \lstinline|False|, and \lstinline|Zero|) immediately
return their corresponding semantic values wrapped in the \lstinline|some| constructor.
For the composite cases, the function uses Lean's do-notation to
sequence computations in the \lstinline|Option| monad. The successor,
predecessor, and zero-test operations (\lstinline|Succ|, \lstinline|Pred|,
and \lstinline|IsZero|) first recursively interpret their subexpressions
with the decremented fuel value, then apply auxiliary semantic functions
(\lstinline|vsucc|, \lstinline|vpred|, and \lstinline|viszero|) to the
resulting values. These functions only compute the resulting value and
wrap them using \lstinline|Option.some| constructor. As an example,
we present the definition of \lstinline|viszero|:
\begin{lstlisting}
def viszero (v : Val .NAT) : Option (Val .BOOL) :=
  match v with
  | .Zero => .some .True
  | _     => .some .False
\end{lstlisting}
which return \lstinline|.True| when the input value is \lstinline|.Zero|
and \lstinline|False|, otherwise. We omit definitions for \lstinline|vsucc| and
\lstinline|vpred| for brevity, since they follow a similar pattern.
Finally, the conditional expression \lstinline|If| evaluates its
boolean condition \lstinline|e1|, then pattern matches on the
resulting value to determine which branch to evaluate.

Type soundness result states that ``well-typed programs cannot
go wrong'' which, in simple terms, means that a well-typed program
cannot result in a runtime error~\cite{Pierce02}.
Following~\cite{Amin17}, we state the theorem by specifying that
the interpreter results are either \lstinline|.none| or a value of
the same type of the input expressions. The Lean 4 encoding of this
theorem is:

\begin{lstlisting}
theorem soundness
  : ∀ fuel t (e : Exp t) r, interp fuel e = r →
    r = .none ∨ ∃ (v : Val t), r = .some v
\end{lstlisting}

We prove this result by induction over the \lstinline{fuel} argument using
the definition of the interpreter. The interested reader could consult the
proof details in the on-line repository~\cite{pegrepository}.

\section{Formalizing Pest operators}~\label{sec:pest-formalization}

% We start by presenting the syntax of the new PEG operators, together with
% some examples. Next, we present the formal semantics of the new operators,
% giving details on repetition indexes are normalized. Finally, we describe
% our extension of Ribeiro et. al. type system~\cite{Ribeiro19} which
% guarantee termination of grammars using the new operators.

\paragraph{New operators syntax}
The Pest tool includes two types of operators: stack and repetition operators.
The syntax of these operators are defined next:
\[
\begin{array}{lcl}
  e & \to  & ...\,\mid\,e^i\,\mid\,e[i,i]\,\mid\,\mathbf{push}\,\mid\,\mathbf{pop}\,\mid\,\mathbf{peek}
             \,\mid\,\mathbf{drop}\\
    & \mid & \mathbf{peekall}\,\mid\,\mathbf{dropall}\\
  i & \to  & \mathbf{top.tonat}\,\mid\,\mathbf{top.length}\,\mid\,i \circ i\,\mid\,v\\
  v & \to  & \mathbf{infty}\,\mid\,n,m\in\mathbb{N}\\
  \sigma & \to & s : \sigma\,\mid\,[] \\
  \circ & \to & + \,\mid\, \times
\end{array}
\]
Meta-variable $s$ denotes an arbitrary string using grammar's alphabet.
We let meta-variable $\sigma$ denote the parser stack and we represent it
using Haskell-style list notation.

\paragraph{Semantics of new operators}
Now that we have introduced the syntax and the informal semantics of the new Pest
operators, we turn our attention to its formal semantics. One important aspect
of Pest repetition operators is the usage of indexes. Before evaluating any
of these new operators we need to reduce indexes to a normal form.
We consider that an index is in normal form if it is a number or the constant
\textbf{infty}. We let meta-variable $v$ denote an arbitrary index normal form.
In simple terms, the index normalization algorithm performs the following steps:

\begin{enumerate}
  \item Evaluate $\mathbf{top.length}$ and $\mathbf{top.tonat}$, if
    they are present in the current index. The usage of these operators
    can result in a run-time error if the parser stack is empty or the
    stack first element can not be converted to a number (for $\mathbf{top.tonat}$).
  \item Reduce all arithmetic operations using their usual meaning and
    the following reduction rules
    when the constant \textbf{infty} is present:
\[
  \begin{array}{c}
    n \circ \mathbf{infty} \rightarrow \mathbf{infty}\\
    \mathbf{infty}\circ n \rightarrow \mathbf{infty}\\
    \mathbf{infty}\circ\mathbf{infty} \rightarrow \mathbf{infty}\\
  \end{array}
\]
where $\circ$ denotes either addition or multiplication and $n$ a
numeric value.
\end{enumerate}

The semantics is denoted by the inductive judgment
(Figures~\ref{fig:newsemantics},~\ref{fig:exactsemantics} and~\ref{fig:intervalsemantics})
$(e, w, \sigma) \Rightarrow_G r$,
where $e$ is the current expression being executed, $w$ is the current input string,
$\sigma$ is the current input parser stack, $G$ is the grammar and $r$ is the output
which can be: 1) success, denoted by triple $(w_1,w_2,\sigma')$ where $w_1$
is the parsed prefix, $w_2$ is the remaining input and $\sigma'$ is the resulting
stack from executing the expression $e$, or; 2) failure, denoted by $\bot$.
The inclusion a stack requires changing the semantics of all standard PEG
operators to thread the current stack state through parser execution.
We omit these modified rules, since they don't add any interesting insight.
\begin{figure*}[h!]
  \[
    \begin{array}{cc}
      \infer[_{\{PUSH_S\}}]
             {(\mathbf{push}\:e,w,\sigma) \Rightarrow_G (w_1, w_2, (w_1:\sigma'))}
             {(e,w,\sigma)\Rightarrow_G (w_1,w_2,\sigma')} &
      \infer[_{\{POP_S\}}]
            {(\mathbf{pop},w,w_1\,:\,\sigma)\Rightarrow_G (w_1,w_2,\sigma)}
            {\exists w_2. w = w_1w_2} \\ \\
        \infer[_{\{DROP_F\}}]
              {(\mathbf{drop},w,[])\Rightarrow_G \bot}{} &
      \infer[_{\{POP_F\}}]
            {(\mathbf{pop},w,\sigma)\Rightarrow_G \bot}
            {\neg \exists w_1\,w_2\, \sigma'. w = w_1w_2 \land \sigma = w_1\,:\,\sigma'}\\ \\
      \multicolumn{2}{c}{
        \infer[_{\{PEEK_F\}}]
              {(\mathbf{peek},w,\sigma)\Rightarrow_G \bot}
              {\neg \exists w_1\,w_2\, \sigma'. w = w_1w_2 \land \sigma = w_1\,:\,\sigma'}
      }\\\\
      \multicolumn{2}{c}{
        \infer[_{\{PEEKALL_S\}}]
              {(\mathbf{peekall},w,\sigma)\Rightarrow_G (w', w'',\sigma)}
              {\sigma = [w_0, ..., w_n] & w' = w_0...w_n & \exists w'' . w = w'w''}
      }\\ \\
      \infer[_{\{PEEK_S\}}]
              {(\mathbf{peek},w,w_1\,:\,\sigma)\Rightarrow_G (w_1,w_2,w_1 \,:\,\sigma)}
              {\exists w_2. w = w_1w_2} &
      \infer[_{\{DROP_S\}}]
            {(\mathbf{drop},w,w_1\,:\,\sigma)\Rightarrow_G (\epsilon,w,\sigma)}
            {}\\ \\
      \multicolumn{2}{c}{
        \infer[_{\{PEEKALL_F\}}]
              {(\mathbf{peekall},w,\sigma)\Rightarrow_G \bot}
              {\sigma = [w_1, ..., w_n] & w' = w_1...w_n & \neg \exists w'' . w = w'w''}
      }\\ \\
         \infer[_{\{PUSH_F\}}]
              {(\mathbf{push}(e),w,\sigma)\Rightarrow_G \bot}
              {(e,w,\sigma)\Rightarrow_G \bot} &
             \infer[_{\{DROPALL_F\}}]
                   {(\mathbf{dropall}, w, [])\Rightarrow_G \bot}
                   {}
          \\ \\
          \multicolumn{2}{c}{
           \infer[_{\{DROPALL_S\}}]
                 {(\mathbf{dropall}, w, w_1 : \sigma)\Rightarrow_G (\epsilon, w, [])}
                 {}
          }
  \end{array}
  \]
  \caption{Semantics of the stack operators}
  \label{fig:newsemantics}
\end{figure*}

Figure~\ref{fig:newsemantics} presents the semantics for the stack manipulation operators.
Rules $PUSH_S$ and $PUSH_F$ formalize the behaviour of the push operator, which pushes the
prefix parsed by the push's expression argument into the parser stack or it fails when
the expression fails. The semantics of the \textbf{pop} operator is defined by rules
$POP_S$ and $POP_F$, which removes the topmost element of the stack and matching it
against the input string. When such a matching is not possible, we have a failure.
The only difference between the semantics for the \textbf{peek} operation and \textbf{pop}
is that the former keeps the top element in the result stack.
The \textbf{drop} operator removes the top element of the stack without matching it
against the current input or it fails when the stack is empty (check rules $DROP_S$ and
$DROP_F$). Rules $PEEKALL_S$ and $PEEKALL_F$ use the whole stack content to match
the current input, failing if all the stack elements do not form a prefix of the current
input and $DROPALL_S$ simply empties the parser stack.

\begin{figure*}[h!]
  \[
    \begin{array}{c}
      \infer[_{\{EX_{S1}\}}]
            {(e^i, w, \sigma)\Rightarrow_G (\epsilon,w,\sigma)}
            {\mathbf{norm}(i,\sigma) = 0} \\ \\
        \infer[_{\{EX_{S2}\}}]
              {(e^i,w,\sigma) \Rightarrow_G (w_1w'_1,w'_2,\sigma')}
              {\begin{array}{c}
                \mathbf{norm}(i,\sigma) = n \\
                (e,w,\sigma)\Rightarrow_G(w_1,w_2,\sigma_1)
               \end{array} &
               \begin{array}{c}
                 0 < n \\
                 (e^{n - 1},w_2,\sigma_1)\Rightarrow_G (w'_1,w'_2,\sigma')
               \end{array}}\\ \\
          \infer[_{\{EX_{F1}\}}]
                {(e^i,w,\sigma)\Rightarrow_G \bot}
                {\begin{array}{c}
                  \mathbf{norm}(i,\sigma) = n \\
                  0 < n \quad (e,w,\sigma)\Rightarrow_G \bot
                \end{array}} \\ \\
            \infer[_{\{EX_{F2}\}}]
                  {(e^i,w,\sigma) \Rightarrow_G \bot}
                  {\begin{array}{c}
                    \mathbf{norm}(i,\sigma) = n \\
                    (e,w,\sigma)\Rightarrow_G(w_1,w_2,\sigma_1)
                  \end{array} &
                  \begin{array}{c}
                    0 < n \\
                    (e^{n - 1},w_2,\sigma_1)\Rightarrow_G \bot
                  \end{array}}
    \end{array}
  \]
  \caption{Semantics for the exact repetition operator}
  \label{fig:exactsemantics}
\end{figure*}

The semantics for the exact repetition operator is presented in
Figure~\ref{fig:exactsemantics}.
The first rule, $EX_{S1}$, shows that for any expression $e$, $e^0$ consumes the
empty string from the input and does not change the parser stack. The second rule,
$EX_{S2}$, shows that to evaluate $e^i$, we first normalize the index $i$ to obtain
a numeric value, $n$. Next, we check if $0 < n$ and parse the input using $e$ producing
a remaining suffix $w_2$, which is used as input for running the expression $e^{n - 1}$.
We say that the repetition operator $e^i$ fails if any iteration of $e$ results in a failure,
as formalized by rules $EX_{F1}$ and $EX_{F2}$.


\begin{figure*}[h!]
  \[
    \begin{array}{cc}
      \infer[_{\{Inter_{1}\}}]
            {(e[i,j],w,\sigma) \Rightarrow_G r}
            {\begin{array}{c}
               \mathbf{norm}(i,\sigma) = n \\
               \mathbf{norm}(j,\sigma) = n \\
               (e^n,w, \sigma)\Rightarrow r
           \end{array}} &
      \infer[_{\{Inter_2\}}]
            {(e[i,j],w,\sigma)\Rightarrow_G (w_1,w_2,\sigma')}
            {\begin{array}{c}
              \mathbf{norm}(i,\sigma) = n \\
              \mathbf{norm}(j,\sigma) = m \\
              n < m \\
              (e^{m}, w, \sigma)\Rightarrow_G (w_1,w_2,\sigma')
          \end{array}} \\ \\
      \infer[_{\{Inter_{F1}\}}]
            {(e[i,j], w, \sigma) \Rightarrow_G \bot}
            {\begin{array}{c}
                \mathbf{norm}(i,\sigma) = n\\
                \mathbf{norm}(j,\sigma) = m\\
                n > m\\
            \end{array}} &
        \infer[_{\{Inter_3\}}]
              {(e[i,j], w, \sigma) \Rightarrow_G r}
              {\begin{array}{c}
                  \mathbf{norm}(i,\sigma) = n \\
                  \mathbf{norm}(j,\sigma) = m \\
                  n < m \\
                  (e^{m},w,\sigma)\Rightarrow_G \bot  \\
                  (e[n,m - 1],w, \sigma)\Rightarrow_G r
              \end{array}}
        \\ \\
        \multicolumn{2}{c}{
          \infer[_{\{Inter_{I}\}}]
                {(e[i,j], w, \sigma) \Rightarrow_G r}
                {\begin{array}{c}
                    \mathbf{norm}(i,\sigma) = n\\
                    \mathbf{norm}(j,\sigma) = \mathbf{infty} \\
                    (e^n\,e^{\star},w,\sigma)\Rightarrow_G r
                  \end{array}
              }}
    \end{array}
  \]
  \centering
  \caption{Semantics for the interval repetition operator}
  \label{fig:intervalsemantics}
\end{figure*}

The last set of semantics rules deals with interval repetition operator,
$e[i,j]$. Rule $Inter_{1}$ parses the input using $e^n$, when both of indexes
normalize to the same numeric value $n$. The rule $Inter_2$ shows how to run
an interval repetition when the normalized indexes are such that $n < m$,
where $\mathbf{norm}(i,\sigma) = n$ and $\mathbf{norm}(j,\sigma) = m$:
it tries to parse the input using $e^m$, if it succeeds then the result is
returned. In case of failure of $e^m$, rule $Inter_3$ tries
to parse the input using the expression $e[n, m -1]$.
Rule $Inter_{F1}$ specifies that an interval repetition fails when the formalized value of $i$
is greater than normalized value of $j$.
The rule deals with the situation when the upper limit bound normalizes to \textbf{infty}.
In this situation, expression $e[i,j]$ is interpreted as $e^n e^*$, where $n$ is the
result of normalizing the index $i$.

\subsection{Type system and its inference algorithm}~\label{sec:typesystem}

Now, we turn our attention to extending a type system to avoid
non termination problems during a grammar execution. We follow
the approach of defining typing rules and its corresponding
type inference using the strategy of constraint generation
followed by a solving step~\cite{Ribeiro19,Cardoso23, Pottier05}.

Before starting discussing the typing rules for the new operators,
we need to extend the $head$ function. In Figure~\ref{fig:new-head-set},
we present the new equations which handle both repetions and stack
operators.

\begin{figure}[H]
  \[
    \begin{array}{lcl}
      head(e^i) & = & head(e)\\
      head(e[i,j]) & = & head(e)\\
      head(\mathbf{push}(e)) & = & head(e)\\
      head(\mathbf{pop}) & = & \emptyset\\
      head(\mathbf{peek}) & = & \emptyset\\
      head(\mathbf{drop}) & = & \emptyset\\
      head(\mathbf{dropall}) & = & \emptyset\\
      head(\mathbf{peekall}) & = & \emptyset\\
    \end{array}
  \]
  \caption{Head-set equations for the new operators}
  \label{fig:new-head-set}
\end{figure}
The first three equations shows that the head-set of
repetitions and the push operator are equal to their
parameter expression set. All the remaining stack operators
have their head-set equal to $\emptyset$, since they do not
refer to non terminals.

The following subsections present the extensions needed to
support typing the new operators.

\subsubsection{Typing rules for the new operators}

We start by defining the typing rules for our new PEG repetition operators.
Figure~\ref{fig:typing-repetition} presents the typing rules for the new
repetition operators. Typing rules make use of the following notations:
1) $\mathbf{norm}^*(i)$ denotes a static analysis
version of the normalization algorithm in which operators \textbf{top.tonat}
and \textbf{top.length} are replaced by 0\footnote{This is necessary because these
operators value depends on the current parser stack state, which is unknown
during type inference.}; 2) we let $\tau^v$ denote the
type $\langle v = 0 \lor \tau.null, \tau.head\rangle$ and
3) notation $\tau[v_1,v_2]$ denote the type
$\langle v_1 = 0 \lor \tau.null, \tau.head\rangle$.


The first type system rule shows that an exact repetition is considered
nullable if either its underlying expression is nullable or its normalized index
is equal to $0$. The head-set for an exact repetition is
the same of its parameter expression. We also restrict that normalized indexes can
only be numeric values, $n$.
Typing of the interval repetition operator is similar to exact repetition, but
we allow the upper interval bound to be \textbf{infty}.
Intervals are considered nullable when
its parameter expression is nullable or either the normalized left index, $n$,
is equal to zero. We also do not allow intervals that have an \textbf{infty}
upper bound to be nullable, since such parsing expression will diverge
whenever they succeed without consuming any input.

\begin{figure}[H]
  \[
    \begin{array}{cc}
      \infer{\Gamma \vdash e^i : \tau^n}
            {\begin{array}{c}
              \mathbf{norm}^*(i) = n\quad \Gamma \vdash e : \tau \\
          \end{array}} &
        \infer{\Gamma\vdash e[i,j] : \tau[n,v] }
              {\begin{array}{ccc}
                  \Gamma \vdash e : \tau &
                  \mathbf{norm}^*(i) = n &
                  \mathbf{norm}^*(j) = v \\
                  \multicolumn{3}{c}{
                  \neg (v \equiv \mathbf{infty} \land\tau.null)}
               \end{array}}
    \end{array}
  \]
  \caption{Typing rules for new PEG repetition operators}
  \label{fig:typing-repetition}
\end{figure}

Figure~\ref{fig:typing-stack} presents the typing rules for the PEG stack
operators. The first rule specifies that the type for \textbf{push}($e$) is
the same as $e$, since it simply pushes the $e$'s consumed input and its
head-set will be the same as $e's$ head-set. Since all other stack operators
can succeed without consuming input, all of them are considered nullable.
As presented in Figure~\ref{fig:new-head-set}, the head-set of \textbf{pop},
\textbf{drop}, \textbf{peek}, \textbf{peekall} and \textbf{dropall} are the
empty set.

\begin{figure}[H]
  \[
    \begin{array}{ccc}
      \infer{\Gamma \vdash \mathbf{push}(e) : \tau}
            {\Gamma \vdash e : \tau}
      &
      \infer{\Gamma \vdash \mathbf{pop} : \langle true, \emptyset\rangle}{} &
      \infer{\Gamma \vdash \mathbf{peek} : \langle true, \emptyset\rangle}{} \\ & \\
      \infer{\Gamma \vdash \mathbf{peekall} : \langle true, \emptyset\rangle}{} &
      \infer{\Gamma \vdash \mathbf{drop} : \langle true, \emptyset\rangle}{} &
      \infer{\Gamma \vdash \mathbf{dropall} : \langle true, \emptyset\rangle}{}
    \end{array}
  \]
  \caption{Typing rules for new PEG stack operators}
  \label{fig:typing-stack}
\end{figure}

\subsection{Extending the type inference algorithm}

In this section we describe the necessary extensions to the type inference algorithm
created by Cardoso et. al.~\cite{Cardoso23} to support the new repetitions and
stack operators.

\paragraph{Extending the constraint syntax} We start by including index normal forms
on the constraint grammar (taken from~\cite{Cardoso23}). The non terminal $v$ denotes
a normal form of an index expression, which can be a numeric value or
constant \textbf{infty}. The complete constraint grammar is as follows:

  \[
    \begin{array}{lcl}
      \tau & \to & \alpha\,\mid\,\langle b, S \rangle\\
      v & \to & n\,\mid\,\mathbf{infty}\\
      \mathbf{t} & \to & A\,\mid\,b\,\mid\,S\,\mid\,\tau\,\mid\,v\\
      C    & \to & \mathbf{true}\,\mid\,\mathbf{false}\,\mid\,\mathbf{t} \equiv \mathbf{t}\,\mid\,C \land C\,\mid\, \exists
                   \alpha.C\,\mid\, \mathbf{def}\:A :\tau\:\mathbf{in}\:C
    \end{array}
  \]
Constraints are first-order logic formulas over types and their meaning
can be defined inductively. We omit the constraint semantics, since its
definition is the same presented in~\cite{Cardoso23}.

\paragraph{Extending constraint generation and solving} After including index expressions
in the constraint syntax, we need to include equations for the new operations
in the constraint generation algorithm.
In order to generate constraints for the exact repetition operator, we use
an existential quantifier to define a new variable, $\alpha$, that is passed
as argument to the generation of constraints for $e$. The algorithm creates
an equality between the argument type, $\tau$, and the repetition operator
type over $\alpha$. We also include an inequality to assert that the
normalized index should not be \textbf{infty}.
\[
  \langle\langle e^i : \tau \rangle\rangle = \exists \alpha . \langle\langle e,\alpha\rangle\rangle \land \tau \equiv \alpha^{\mathbf{norm}^*(i)} \land \mathbf{norm}^*(i) \not\equiv \mathbf{infty}\\
\]
Next, we consider the generation of constraints for the interval repetition operator, which
follows the same pattern of the exact repetition: we generate the constraint for sub-expression
$e$, considering an existential quantified type variable and create an equality between the variable
type and the input argument type, $\tau$. We also impose that the interval lower-bound
cannot be \textbf{infty} and that if the upper-bound is \textbf{infty}, the interval
underlying expression cannot be nullable since it could make the parser loop indefinitely.
\[
  \begin{array}{rl}
    \langle\langle e[i,j] : \tau \rangle\rangle =   & \exists \alpha . \langle\langle e, \alpha\rangle\rangle \land \tau \equiv \alpha[\mathbf{norm}^*(i),\mathbf{norm}^*(j)]\\
                                            \land{} & \mathbf{norm}^*(i) \not\equiv\mathbf{infty} \\
                                            \land{} & \neg (\mathbf{norm}^*(j) \equiv \mathbf{infty} \land \alpha.null)\\
  \end{array}
\]
The constraint for the \textbf{push} operator is just the constraint of its argument expression and
all other stack expressions just generate an equality between its input argument type and
$\langle true,\emptyset\rangle$.

\begin{figure}[H]
  \[
    \begin{array}{lcl}
      \langle\langle \mathbf{push}(e) : \tau\rangle\rangle & = & \langle\langle e : \tau \rangle\rangle \\
      \langle\langle \mathbf{pop} : \tau \rangle\rangle & = & \tau \equiv \langle true, \emptyset \rangle\\
      \langle\langle \mathbf{drop} : \tau \rangle\rangle & = & \tau \equiv \langle true, \emptyset \rangle\\
      \langle\langle \mathbf{dropall} : \tau \rangle\rangle & = & \tau \equiv \langle true, \emptyset \rangle\\
      \langle\langle \mathbf{peek} : \tau \rangle\rangle & = & \tau \equiv \langle true, \emptyset \rangle\\
      \langle\langle \mathbf{peekall} : \tau \rangle\rangle & = & \tau \equiv \langle true, \emptyset \rangle\\
    \end{array}
  \]
  \caption{Extending the constraint generation algorithm}
  \label{fig:constraint-gen}
\end{figure}

The inclusion of index expressions do not change the constraint solver specification~\cite{Cardoso23}.
Normalized indexes are just integers or the \textbf{infty} and they only appear as part of the
nullability field of a type, which has its dedicated equation in solver specification, or as part
of an inequality, which can be easily solved by an SMT solver or a simple custom solving algorithm.

Since the proposed changes (the inclusion of indexes and new operators) do not change the
constraint solver, just the constraint generation algorithm, we argue that the
soundness and completeness argument~\cite{Cardoso23} for the type inference algorithm
also holds for our Pest-style grammars. This lead to our main result:

\begin{Theorem}[Termination of typed grammars]
  Let $G = (V,\Sigma,R, e_s)$ be a well-typed grammar which may use any of Pest
  repetition or stack operators. Then, for all $w$, exists $r$
  such that $(e_s,w,[])\Rightarrow r$.
\end{Theorem}

The only way for a parsing expression to loop over an input is due to
the presence of left-recusive rules or nullable expressions
under a star operator in a grammar~\cite{Ford2004}. Since the type system (and its
inference algorithm) reject any grammar that have these issues,
any well-typed grammar is guaranteed to terminate.

\section{Formalizing the type soundness theorem in Lean 4}\label{sec:typesoundnesslean}

In the type theory literature, the Wright and Feilleisen approach for type soundness is
dominant~\cite{Wright94}. The idea is to split the proof of type soundness in two lemmas:
1) the preservation lemma, which states that an execution step on a well-typed program
preserves types; and 2) the progress lemma, which says that every well-typed program is
either a final result of a computation (a value) or it can give one more execution step.
As argued by Amin and Rompf~\cite{Amin17}, since such proof style is applicable only on
the context of languages which do have a small-step / reduction semantics, they propose the use of
definitional interpreters to formalize soundness in modern proof assistants.
In this section, we describe our Lean 4 formalization of the type soundness results for our
type system for Pest grammars using an intrisincally-typed approach, as argued by~\cite{plfa22.08}.
We start by describing some dependently-typed toolkit
used in the formalization in Section~\ref{sec:depkit}. Indexes and their
evaluation procedures are described in Section~\ref{sec:index-lean}. Next, we present
the intrisincally-typed encoding of grammars (Section~\ref{sec:pest-syntax}) and then,
in Section~\ref{sec:pest-interpreter}, we describe our implementation of a definitional
interpreter for grammars and we explain how it enables a simple formalization of the type
soundness theorem for grammars.

\subsection{Dependently-typed utilities}~\label{sec:depkit}

In the development of our intepreter, we decided to drop variable names and represent
them as DeBruijn indexes~\cite{DeBruijn72}. There are different approaches for implementing
De Bruijn indexes in proof assistants~\cite{Ni24}, but in this work we choose an
inductive type that represent the set of the first $n$ natural numbers.
\begin{lstlisting}
inductive fin : ℕ -> Type where
| zero : ∀ {n}, fin (n + 1)
| succ : ∀ {n}, fin n → fin (n + 1)
\end{lstlisting}
Type \lstinline|fin n| holds exactly $n$ inhabitants which make it perfect to represent
the finite set of variables of a grammar. Another useful type which is the length
indexed list type, usually named as vector:
\begin{lstlisting}
inductive vec (A : Type) : ℕ → Type where
| nil : vec A 0
| cons : ∀ {n}, A → vec A n → vec A (n + 1)
\end{lstlisting}
Type \lstinline|vec A n| denotes $n$ element lists of type \lstinline|A|. Constructor
\lstinline|nil| builds an empty list, which has size $0$ and \lstinline|cons| inserts
a new element into a $n$ element list to create a new list containing $n + 1$ values.
The \lstinline|vec.lookup| function provides a safe way to access elements in a
length-indexed vector by using a bounded index. It takes a vector \lstinline|v|
which stores \lstinline|n| elements of type \lstinline|A| and an index
\lstinline|i : fin n| (which is guaranteed to be less than n), and returns the
element at vector's position \lstinline|i|. The function works recursively
through pattern matching: when the index is \lstinline|fin.zero|, it returns the
head of the vector; when the index is \lstinline|fin.succ i1|, it recursively
looks up \lstinline|i1| in the tail of the vector \lstinline|v1|.
This implementation eliminates the possibility of index-out-of-bounds errors
at compile time, since only valid vector positions can be passed as arguments to
this function.
\begin{lstlisting}
def vec.lookup {n}{A : Type}(v : vec A n)(i : fin n) : A :=
  match i, v with
  | fin.zero, vec.cons x _ => x
  | fin.succ i1, vec.cons _ v1 =>
    vec.lookup v1 i1
\end{lstlisting}
The last piece of our toolkit is the \lstinline|vec.all| type. It is a dependent
predicate that expresses that every element in a vector satisfies a given
property \lstinline|P|. It's defined inductively to mirror the structure of
vectors: the nil constructor states that the empty vector trivially satisfies
\lstinline|vec.all P|, while the \lstinline|cons| constructor requires both
a proof that the head element \lstinline|x| satisfies \lstinline|P x| and a
proof that the tail \lstinline|v| satisfies \lstinline|vec.all P v|.
\begin{lstlisting}
inductive vec.all {A : Type}(P : A → Type) : ∀ {n}, vec A n → Type where
| nil : vec.all P nil
| cons : ∀ {n}{x : A}{v : vec A n}, P x → vec.all P v → vec.all P (vec.cons x v)
\end{lstlisting}
Finally, function \lstinline|all.lookup| allow us to get a proof that a
vector element satisfies a property \lstinline|P| from an index \lstinline|i|.
\begin{lstlisting}
def all.lookup {A : Type}{P : A → Type}{n}{v : vec A n}(xs : vec.all P v)(i : fin n)
  : P (vec.lookup v i) :=
  match i, xs with
  | fin.zero, vec.all.cons x _ => x
  | fin.succ i1, vec.all.cons _ v1 =>
    all.lookup v1 i1
\end{lstlisting}
The implementation works by structural recursion, matching on both the index
and the proof: when \lstinline|i| is \lstinline|fin.zero|, it extracts the
proof for the head element; when \lstinline|i| is \lstinline|fin.succ i1|, it
recursively looks up the proof in the tail.

\subsection{Indexes}\label{sec:index-lean}

We start by representing the syntax of indexes as the following
inductive type:
\begin{lstlisting}
inductive Index : Type where
| nat : ℕ → Index
| infty : Index
| tonat : Index
| length : Index
| add : Index → Index → Index
| mult : Index → Index → Index
\end{lstlisting}
Type \lstinline|Index| has a constructor for each index syntatic form: constants
are represented by the constructor \lstinline|nat|, the infinite index is denoted
by \lstinline|infty|, stack query indexes \lstinline|tonat| and \lstinline|length|
denote the top stack element corresponding natural number value or its length, respectively.
Addition and multiplication of indexes are built using constructors \lstinline|add| and
\lstinline|mult|. The result of evaluating indexes are represented by values,
which are either numeric
constant or the special value \lstinline|infty|, which are denoted by the \lstinline|Value|
type.
\begin{lstlisting}
inductive Value : Type where
| nat : ℕ → Value
| infty : Value
\end{lstlisting}
In order to reduce indexes, we use two distinct functions: the first, \lstinline|norm_static|
is the implementation of $\mathbf{norm}^*$ normalization procedure used by type system.
This definition proceeds by pattern matching over the index structure and provides no
surprises, except that stack query indexes will always return 0, since we do not have access
to the parsing stack during type checking and the need to propagate the \lstinline|infty|
value through recursive calls of the normalization of arithmetic operation arguments.
\begin{lstlisting}
def norm_static (i : Index) : Value :=
  match i with
  | .nat n => .nat n
  | .infty => .infty
  | .tonat => .nat 0
  | .length => .nat 0
  | .add i1 i2 =>
    match norm_static i1, norm_static i2 with
    | .nat n1, .nat n2 => .nat (n1 + n2)
    | .infty , _ => .infty
    | _, .infty => .infty
  | .mult i1 i2 =>
    match norm_static i1, norm_static i2 with
    | .nat n1, .nat n2 => .nat (n1 * n2)
    | .infty , _ => .infty
    | _, .infty => .infty
\end{lstlisting}
The normalization of indexes used by the semantics is implemented by function \lstinline|norm|,
which takes the current parser stack as and additional argument, together with the
index to be evaluated and return an optional value, since evaluating stack query indexes can
cause runtime errors. Again, most of the \lstinline|norm| is similar to \lstinline|norm_static|.
Differences arise due to dealing with stack querying operations and the additional propagation
of runtime errors (encoded by the \lstinline|Option.none| value) through recursive calls.
\begin{lstlisting}
def norm (i : Index)(s : stack) : Option Value :=
  match i, s with
  | .nat n, _ => .some (.nat n)
  | .infty, _ => .some .infty
  | .tonat, [] => .none
  | .tonat, (s :: _) =>
    match (String.ofList s).toNat? with
    | .none => .none
    | .some i => .some (.nat i)
  | .length, [] => .none
  | .length, (s :: _) => .some (.nat s.length)
  | .add i1 i2, stk =>
    match norm i1 stk, norm i2 stk with
    | .some (.nat n1), .some (.nat n2) =>
        .some (.nat (n1 + n2))
    | .some .infty, _ => .some .infty
    | _ , .some .infty => .some .infty
    | _, _ => .none
  | .mult i1 i2, stk =>
    match norm i1 stk, norm i2 stk with
    | .some (.nat n1), .some (.nat n2) =>
        .some (.nat (n1 * n2))
    | .some .infty, _ => .some .infty
    | _ , .some .infty => .some .infty
    | _, _ => .none
\end{lstlisting}
The definitions of \lstinline|norm| and \lstinline|norm_static| are related by some
lemmas, which are proved by a straightforward induction over the input index. We use
such lemmas to get a contradiction in the interpreter definition, to satisfy Lean's
totality checker which is not able to determine that such situations are impossible.

\subsection{Intrisincally-typed syntax for grammars}~\label{sec:pest-syntax}

In this section we describe our design decisions to encode our proposed type system
as an intrisincally-typed syntax definition using Lean 4.
We start by defining the representation of types, which mimics the exact definition
of the type system: a record formed by a boolean and the headset. Since we are representing
variables as DeBruijn indexes, which are values of the \lstinline|fin| type, the headset is
just a list of \lstinline|fin n| values, where \lstinline|n| is the number of variables in
the grammar. Typing contexts are just vectors of types.
\begin{lstlisting}
abbrev headset (n : ℕ) := List (fin n)

structure type (n : ℕ) : Type where
  nullable : Bool
  headset : headset n

abbrev ctx (n : ℕ) := vec (type n) n
\end{lstlisting}
Operations over types are implemented as functions to combine types. As an example,
the sum of two types is realized by \lstinline|type.sum| function. Operations for
product, negation, star and repetition operators follow a similar structure and
are ommited for brevity.
\begin{lstlisting}
def type.sum {n}(t1 t2 : type n) : type n :=
  type.mk (t1.nullable || t2.nullable)
          (t1.headset ++ t2.headset)
\end{lstlisting}
We use the typing rules from Figures~\ref{fig:typesystem},~\ref{fig:typing-repetition}
and \ref{fig:typing-stack} as the start point from our implementation. First, we
define a predicate to establish when a variable reference is valid:
\begin{lstlisting}
def validvar {n}(sig : ctx n)(v : fin n) : Prop :=
  let t : type n := vec.lookup sig v
  match t with
  | type.mk _ h => v ∉ h
\end{lstlisting}
This definition encodes the premise present in the typing rule for variables in
Figure~\ref{fig:typesystem}, which requires that $A \not\in \tau.head$, when
typing a variable $A$ with type $\tau$. The predicate \lstinline|validvar sig v|
holds when the variable \lstinline|v| does not appear in its own head-set,
thereby preventing left recursion in the grammar.

Now, we have the necessary ingredients to develop the syntax of parsing expressions
that only allow the construction of well-typed grammars. The Lean 4 representation of
this type is as follows:
\begin{lstlisting}
inductive pexp {n}(sig : ctx n) : type n -> Type where
| epsilon : pexp sig (type.mk true [])
| chr : Char → pexp sig (type.mk false [])
| var : (v : fin n) → validvar sig v → pexp sig (vec.lookup sig v)
| cat : ∀ {t1 t2}, pexp sig t1 → pexp sig t2 → pexp sig (type.prod t1 t2)
| choice : ∀ {t1 t2}, pexp sig t1 → pexp sig t2 → pexp sig (type.sum t1 t2)
| not : ∀ {b s}, pexp sig (type.mk b s) → pexp sig (type.mk true s)
| star : ∀ {s}, pexp sig (type.mk false s) → pexp sig (type.mk true s)
| push : ∀ {t}, pexp sig t → pexp sig t
| pop : pexp sig (type.mk true [])
| peek : pexp sig (type.mk true [])
| drop : pexp sig (type.mk true [])
| peekall : pexp sig (type.mk true [])
| dropall : pexp sig (type.mk true [])
| rep : ∀ {n t} i, norm_static i = Value.nat n → pexp sig t → pexp sig (type.tyrep n t)
| bound : ∀ {n v} t i1 i2, norm_static i1 = Value.nat n → norm_static i2 = v →
                           pexp sig t → ¬ (v = .infty ∧ t.nullable = true) →
                           pexp sig (type.mk (Nat.beq n 0 && t.nullable) t.headset)
\end{lstlisting}
Type \lstinline|pexp sig t| denotes a parsing expression with type \lstinline|t|
in a typing context \lstinline|sig|. We start describing the constructors encoding
the typing rules for standard PEG operators. The constructors \lstinline|epsilon|
and \lstinline|chr| shows that expressions $\epsilon$ and $a \in \Sigma$
have types $\langle true, \emptyset \rangle$ and $\langle false, \emptyset\rangle$,
respectively. The constructor for variables uses the predicate \lstinline|validvar| to
restrict the variable to be not contained on its headset. Concatenation and choice
simply use the operations \lstinline|type.prod| and \lstinline|type.sum| to combine
the types of its subexpressions. Constructor \lstinline|not| denotes the PEG not
predicate and it simply forces the type of its subexpression to be nullable, while
keeping the same headset. The constructor for the star operator forces that its
subexpression should not be nullable by requiring it to have the type
\lstinline|pexp sig (type.mk false s)|, for some headset s.

Next, we turn our attention to stack operators. Constructor \lstinline|push| specifies
that expression \lstinline|push e| has the same type as \lstinline|e|. All other
stack operators are considered nullable and having an empty head-set. The constructor
for exact repetition operator, \lstinline|rep|, demands that its sub-expression is
well-typed and that its index normalizes into a numeric constant. Finally, constructor
\lstinline|bound| express the restrictions over bounded repetition expressions:
1) the sub-expression should be well-typed; 2) the lower bound should always reduce
to a numeric constant and 3) Let \lstinline|v| be the normalized index for the upper-bound.
Then if the sub-expression is nullable, \lstinline|v| should be different from
\lstinline|infty|.

Following~\cite{Ribeiro19}, since we represented variables as De Bruijn indexes, we need to
relate an expression with its type on typing context. Datatype
\lstinline|rules| stores the set of rules (parsing expressions) of a grammar in
a length-indexed list (vector). We use type \lstinline|vec.all| to enforce the relation
between an expression and its type on context.
\begin{lstlisting}
structure rules {n}(sig : ctx n) : Type where
  exprs : vec.all (λ t => pexp sig t) sig
\end{lstlisting}
Finally, we represent a grammar by a record containing a typing context for
its non-terminals types, its rules and the grammar start expression together with
its type.
\begin{lstlisting}
structure grammar : Type where
  n : ℕ
  sig : ctx n
  pegs : rules sig
  starttype : Ty
  start : pexp sig starttype
\end{lstlisting}

\subsection{Definitional interpreter for grammars and its soundness theorem}~\label{sec:pest-interpreter}

In this section we detail the implementation of the interpreter for the
intrinsically-typed syntax for grammars and how it can be used to produce a proof of the
soundness theorem for our type system.
We follow the same strategy as~\cite{Ribeiro19, Amin17} and use a ``fuel'' parameter
to limit the number of recursive calls made by the interpreter. We also need
a special type to denote valid parsing results:
\begin{lstlisting}
inductive Result : Type where
| Ok : word → word → stack → Result
| ParserError : Result
\end{lstlisting}
Type \lstinline|Result| encodes possible PEG execution results: it will either fail
on the current input (constructor \lstinline|ParserError|) or it will return a
pair formed by the consumed prefix, the remaining input and the current stack
state, as represented by the constructor \lstinline|Ok|.

The main intepreter function, \lstinline|parse|, has the following signature:
\begin{lstlisting}
def parse {n}{sig : ctx n}{t : type n}(fuel : ℕ)(rs : rules sig)
          (e : pexp sig t)
          (s : word)(st : stack) : Option Result :=
    ... -- definition of the interpreter cases
\end{lstlisting}
and it is defined by recursion on the structure of the expression \lstinline|e|
and the current fuel value. Whenever the interpreter runs out of fuel, we return
\lstinline|Option.none| to denote such out of gas error.
\begin{lstlisting}
match fuel with
| 0 => .none
| fuel1 + 1 => -- matching over e..
\end{lstlisting}
Next, we consider the semantics for $\epsilon$. It will always succeed by
returning an empty parsed prefix, the current input and the current
stack state as the parsing result.
\begin{lstlisting}
match e with
| pexp.epsilon => pure (.Ok [] s st)
\end{lstlisting}
The parsing of a single character expression needs to check if
it is the first input symbol. Otherwise, we have a failure.
We start by doing a case analysis on the input, returning a
failure when it is empty. For non-empty strings, we check if
the first symbol is equal to the current parsing expression
using function \lstinline|decEq|, which is the decidable
equality test defined in Lean standard library. Whenever the
two characters aren't equal, we return a parse error or a
triple containing the parsed symbol, the remaining input and
the current stack.
\begin{lstlisting}
| pexp.chr c =>
  match s with
  | [] => pure .ParserError
  | (c1 :: s1) =>
    match decEq c c1 with
    | .isTrue rfl => pure (.Ok [c] s1 st)
    | .isFalse _ => pure .ParserError
\end{lstlisting}
For the variable case, we simply execute its right-hand side by
looking up its defining expression in the grammar rules and calling
the intepreter recursively over it.
\begin{lstlisting}
| pexp.var v _ =>
    match all.lookup rs.exprs v with
    | e1 => parse fuel1 rs e1 s st
\end{lstlisting}
Evaluation of the concatenation of two expressions proceeds as follows:
we first evaluate \lstinline|e1| over the input \lstinline|s|, if it results
in a failure, then the whole concatenation fails. If it succeeds, then it
should produce a parsed prefix, a remaining suffix and a result stack. Then,
we execute \lstinline|e2| over the suffix. Again, if it fails, concatenation
fails. When \lstinline|e2| succeeds, the result is the concatenation of
the parsed prefix of its sub-expressions, the remaining input is what
\lstinline|e2| leaves from the input and the result stack will be the value
produced by the execution of \lstinline|e2|. This semantics is expressed by
the following Lean equations.
\begin{lstlisting}
| pexp.cat e1 e2 =>
    match parse fuel1 rs e1 s st with
    | .some (.Ok pre1 suf1 st1) =>
      match parse fuel1 rs e2 suf1 st1 with
      | .some (.Ok pre2 suf2 st2) =>
        pure (.Ok (pre1 ++ pre2) suf2 st2)
      | .some .ParserError => pure .ParserError
      | .none => .none
    | .some .ParserError => pure .ParserError
    | .none => .none
\end{lstlisting}
Next, we consider the semantics for the choice operator: it parses the input
using the first subexpression, if it succeeds then the whole choice expression
will also succeed. If we reach a not enough gas error
(resulting a \lstinline|none| constructor), we simply propagate the error.
When \lstinline|e1| execution results in a parser error, we try to run
\lstinline|e2| over the original input \lstinline|s|.
\begin{lstlisting}
| pexp.choice e1 e2 =>
    match parse fuel1 rs e1 s st with
    | .some (.Ok pre1 suf1 st1) => .some (.Ok pre1 suf1 st1)
    | .some .ParserError => parse fuel1 rs e2 s st
    | .none => .none
\end{lstlisting}
The semantics of the not operator is defined as follows: it runs
its subexpression against the input. A successful result from the
subexpression causes the not operator to produce a parser error.
A failure from the subexpression causes the not operator to succeed,
without consuming input.
\begin{lstlisting}
| pexp.not e1 =>
  match parse fuel1 rs e1 s st with
  | .some (.Ok _ _ _) => pure .ParserError
  | .some .ParserError => pure (.Ok [] s st)
  | .none => .none
\end{lstlisting}
The semantics of the Kleene operator work as follows: we try to parse
the input using the Kleene sub-expression. When it fails, the whole
Kleene succeeds without consuming input. If the sub-expression succeeds,
the interpreter runs the Kleene expression again over the remaining suffix.
The only situation which results in failure is when we ran out of fuel. In this
situation, we just propagate the value \lstinline|.none|.
\begin{lstlisting}

| pexp.star e1 =>
  match parse fuel1 rs e1 s st with
  | .some .ParserError => pure (.Ok [] s st)
  | .some (.Ok pref1 suf1 st1 ) =>
    match parse fuel1 rs (pexp.star e1) suf1 st1 with
    | .some (.Ok pref2 suf2 st2) =>
      pure (.Ok (pref1 ++ pref2) suf2 st2)
    | .some .ParserError =>
      pure (.Ok pref1 suf1 st1)
    | .none => .none
  | .none => .none
\end{lstlisting}
Now, let's focus on the stack operators. We start with \lstinline|push e1|, which
parse the input using the sub-expression \lstinline|e1|. The resulting parsed prefix
is then pushed in the result stack.
\begin{lstlisting}
| pexp.push e1 =>
    match parse fuel1 rs e1 s st with
    | .some (.Ok pref1 suf1 st1) =>
      pure (.Ok pref1 suf1 (pref1 :: st1))
    | .some .ParserError => pure .ParserError
    | .none => .none
\end{lstlisting}
Next, we consider the pop operator, which tries to parse the stack top in the current
input. A parser error is returned whenever the stack is empty or the top string does
not match with the prefix of the current input.
\begin{lstlisting}
| pexp.pop =>
  match st with
  | [] => pure .ParserError
  | (w :: st1) =>
    let n  := List.length w
    let w1 := List.take n s
    if w = w1 then pure (.Ok w  (List.drop n s) st1)
    else pure .ParserError
\end{lstlisting}
Operators \lstinline|peek|, \lstinline|drop|, \lstinline|peekall| and \lstinline|dropall|
have a trivial implementation and are omitted for brevity. Finally, let's focus on
the new repetition operators. First, we show the equations that define
the semantics of the exact repetition operator.
\begin{lstlisting}
| pexp.rep i eq e =>
  match Hv : norm i st with
  | .some (Value.nat 0) => pure (.Ok [] s st)
  | .some (Value.nat (m + 1)) =>
    match parse fuel1 rs e s st with
    | .some (.Ok pref1 suf1 st1) =>
      let i1 := Index.nat m
      let eq1 : norm_static i1 = Value.nat m := by rfl
      match parse fuel1 rs (pexp.rep i1 eq1 e) suf1 st1 with
      | .some (.Ok pref2 suf2 st2) =>
        pure (.Ok (pref1 ++ pref2) suf2 st2)
      | .some .ParserError => pure .ParserError
      | .none => .none
    | .some .ParserError =>
      pure .ParserError
    | .none => .none
\end{lstlisting}
The first step is to obtain the normal form \lstinline|v| of the index of the
exact repetition. When the normal form is a numeric value, then we simply repeat
the parsing of the expression \lstinline|e| over the input. Parser errors and
out of gas are simply propagated in the call stack of the semantics. When the
normal form \lstinline|v| is \lstinline|infty|, we have a contradiction, since
well-typed exact repetition expressions cannot have indexes which normalize to
\lstinline|infty|. Our semantics deal with this contradictory case using a
simple tactic script to discharge it. The reader interested in the details of
this script can check our formalization source code.
The final operator considered in our semantics is the bounded repetition. Its
implementation starts by normalizing their two indexes producing values
\lstinline|v1| and \lstinline|v2|. Then, we consider the following possibilities:
1) If \lstinline|v1 < v2|, we simply execute $e[v1, v2 - 1]$ over the current input;
2) If \lstinline|v1 = v2|, we then parse the input using $e^{v1}$;
3) If \lstinline|v2 = infty|, we parse the input using $e^{v1}e^\star$.

\begin{lstlisting}
| pexp.bound t i1 i2 eq1 eq2 e neq =>
  match Hv1 : norm i1 st, Hv2 : norm i2 st with
  | .some (Value.nat n1), .some (Value.nat n2) =>
    match Ord.compare n1 n2 with
    | Ordering.lt =>
      let i3 := Index.nat n2
      match parse fuel1 rs (pexp.rep i3 (by rfl) e) s st with
      | .none => .none
      | .some .ParserError =>
        let i4 := Index.nat n1
        let i5 := Index.nat (n2 - 1)
        parse fuel1 rs (pexp.bound t i4 i5 (by rfl) (by rfl) e
                                           (by intros H3 ; simp [norm_static] at *)) s st
      | res => res
    | Ordering.eq =>
      let i3 := Index.nat n1
      parse fuel1 rs (pexp.rep i3 (by rfl) e) s st
    | Ordering.gt => .some .ParserError
  | .some (Value.nat n1), .some .infty =>
    if Ht : t.nullable then by
    -- contradictory case.
    -- tactic script omitted for brevity
    else
      let i3 := Index.nat n1
      let e1 : pexp sig (type.mk false t.headset) := by
        have Hn : t.nullable = false := by aesop
        have Heq : t = type.mk t.nullable t.headset := by
          rcases t
          aesop
        rw [Heq] at e
        rw [Hn] at e
        exact e
      parse fuel1 rs (.cat (.rep i3 (by rfl) e) (.star e1)) s st
\end{lstlisting}
The final piece of our semantics is the function that executes a PEG
over an input. Function \lstinline|run| takes a fuel argument, a grammar
and an input string and calls function \lstinline|parse| passing the grammar
rules, its start expression, the input string and an empty stack.
\begin{lstlisting}
def run (fuel : ℕ)(g : grammar)(s : word) : Option Result :=
  parse fuel g.pegs g.start s []
\end{lstlisting}

From the implementation of the interpeter, we can see that the only situation
that it returns \lstinline|.none| is when it runs out of gas. There are two reasons
for exausting the fuel counter: either its initial value was less than the minimum
necessary for the complete parsing of the input or the grammar looped. Since our
type system rejects any grammar with left-recursive rules and expressions which
contain a nullable sub-expression of a Kleene star, grammars that can be executed
using the intepreter cannot loop over the input. So, the only reason for
a \lstinline|.none| result is that the initial fuel value was not enough for
the current grammar and input.

In what follows, we state the soundness theorem. Its proof proceeds by
by induction over the fuel counter and case analysis on the
intrinsically-typed representation of expressions. We omit the tactic
script for brevity. The interested reader could check its details on-line~\cite{pegrepository}.
\begin{lstlisting}
theorem soundness
  : ∀ {fuel}{n}{sig : ctx n}{t}(rs : rules sig)(e : pexp sig t) s st r,
  parse fuel rs e s st = .some r →
  r = Result.ParserError ∨ ∃ pre, ∃ suf, ∃ st, r = .Ok pre suf st
\end{lstlisting}

\section{Examples}~\label{sec:examples}

We implemented a Racket library, called \emph{typed-peg},
which implements the type inference and the semantics of grammars with the proposed new
operators. In this section, we present a complete example grammar and discuss how
our approach avoid the looping behavior reported by Pest users~\cite{PESTBUG}.

\paragraph{A grammar for the Heartbeat packages} One of the most famous binary data format
bug is the so-called ``heartbleed''. The ASN.1 binary interface uses a ``heartbeat''
package to check if the other remote party is still alive. A well-formed heartbeat package
is formed by the following components:
\begin{itemize}
  \item A type field, which indicates what kind of ``heartbeat'' is being used,
  \item A field for the length $l$ --- a two byte value,
  \item A challenge $C$ of length $|C| = l$, and
  \item An arbitrary number of padding bytes.
\end{itemize}

\begin{verbatim}
#lang typed-peg
bit <-- '0' / '1'
byte <-- bit^8
type <-- byte
length <-- push(byte^2)
challenge <-- byte^top.tonat
padding <-- byte*
start: type ~ length ~ challenge ~ padding
\end{verbatim}
The previous grammar is a simple encoding of the heartbeat package format.
A key issue regarding implementing parsers for this format is to ensure that
challenge has the size specified by the length field. Using the stack operator
\textbf{push} we insert the parsed length into the execution stack and retrieve
its numeric value using the index expression \textbf{top.tonat}, which parses
the challenge using the exact repetition operator.

\paragraph{A grammar for the PNG format}

A popular image file format is the Portable Network Graphics (PNG), which
is represented by the following pieces of data: 1)
A format signature (8 bytes), which always corresponds to the
following integer values: 137, 80, 78, 71, 13, 10, 26, 10;
2) A 4 bytes natural number $n$, which represents the data length;
3) 4 bytes chunk type value;
4) The image data, a sequence of $n$ bytes and
5) a 4 bytes CRC checksum.
A possible grammar to specify the PNG format is:
\begin{verbatim}
#lang typed-peg
bit <-- '0' / '1'
byte <-- bit^8
fist-four <-- '137' ~ '80' ~ '78' ~ '71'
snd-four <-- '13' ~ '10' ~ '26' ~ '10'
signature <-- first-four ~ snd-four
length <-- push(byte^4)
type <-- byte^4
data <-- byte^top.tonat
crc <-- byte^4
start: signature ~ length ~ type ~ data ~ crc
\end{verbatim}
Notice that the use of an exact repetition operator combined with
index \textbf{top.tonat} allows us to guarantee that the data block
has the correct size of the format length field.

\paragraph{Rejecting reported looping examples} Now, let's turn our
attention to some expressions which Pest users~\cite{PESTBUG} reported as
non terminating.

\begin{verbatim}
forever1 <-- (push('') ~ pop)*
forever2 <-- popall ~ (popall)*
forever3 <-- forever3
\end{verbatim}
In the first two examples, we have nullable expressions under a star
operator. In the first rule, we have a repetition of \textbf{push} operator, which
has an empty string expression as its argument, followed by \textbf{pop}.
A \textbf{push} is nullable only if its argument is nullable. Since it
has an empty string as argument, the expression \verb|push('')|
accepts the empty string. Next, we have a pop operator which is nullable,
by the definition. In this way, the type system rejects \verb|forever1|,
since it has a nullable expression under a star. The second example is also
ill-typed since it has a \textbf{popall}, a nullable operator, as
argument to a PEG repetition.
The third example shows a simple left recursive rule
which is also rejected by type inference algorithm.

\paragraph{A grammar for an identation sensitive language}

Parsing identation sensitive languages poses a challenge to language developers.
The next grammar is an example of how stack operators can be used to define
a language which uses different levels of identation. The idea is to use the
stack to control the identation level.
\begin{verbatim}
#lang typed-peg
new_line <-- '\n'
char <-- 'a' / ... / 'z' / '0' / ... / '9'
ident_level <-- ' '+ / '\t'+
block_name <-- char+
block <-- block_name ~ new_line ~ new_block*
new_block <-- peekall ~
              push (ident_level) ~
              block ~
              (peekall ~ block)* ~
              drop
start: new_line* ~ block* ~ new_line*
\end{verbatim}
At each block, we start by matching the current block identation
using the operator \textbf{peekall}. Next, we push the possible
next identation level and process a possibly empty set of blocks.
We end a block by removing its identation from the parser stack
using \textbf{drop}.

\paragraph{Discussion}
Our work uses the framework of operational semantics and type
theory to provide a solid foundation to Pest stack and repetition
operators. The main issue pointed by Pest users~\cite{PESTBUG} is failure of
detecting non termination of grammars. The proposed type inference
algorithm can avoid this problem by ruling out any grammar which
has direct or indirect left-recursive rules and nullable expressions
as a parameter to the Kleene star operator.

During our formalization, we notice that the inclusion
of indexes that allow the inspection of the topmost stack item
would increase the tool expressivity, since they allow concise
specifications for some context-sensitive formats.
However, the proposed
typing algorithm does not avoid some kinds of run-time errors, notably
errors caused by \textbf{top.tonat} and \textbf{top.length}, when
the stack is empty or if its top element is not a
string formed only by digits, in the case of \textbf{top.tonat}.
Current version of the Pest tool do not provide support to \textbf{top.nat}
and \textbf{top.length} indexes.

A limitation of our current prototype is that it does not handle different
input encodings. The semantics of index expression \textbf{top.tonat} is
dependent of current input encoding to convert the topmost stack value to
a natural number, if possible. The proper handling of multiple encodings
is left for future work.

While our approach is able to detect non terminating grammars it does not avoid
some runtime errors, notably caused by \textbf{top.tonat} and \textbf{top.length},
since if the parser stack is empty (for both operations) or if its top element is
not a valid number (for \textbf{top.tonat}).

\section{Related work}~\label{sec:related}

The works of Ribeiro et al. ~\cite{Ribeiro19} and Cardoso et. al.~\cite{Cardoso23}
proposed a type system and a type inference algorithm for PEGs that is sound with
respect to Ford's well-formedness predicate. Authors argue that the use
of a type system provide a more predictable and compositional behavior.
Another work that uses types to ensure parsing termination was proposed by
Krishnaswami et. al.~\cite{Krishnaswami19}. The authors define a type system
for $\mu$-regular expressions, an algebraic presentation of context-free
languages, and proved that their type system only accepts LL(1) grammars.
Our work extend both the type system for PEGs and its inference algorithm to
Pest operators while ensuring that no well-typed grammar
will loop over inputs.

The development of formalisms and tools for correct parsing of complex data
formats is subject of some recent works.  Zhang et. al.~\cite{Zhang23} proposed
interval parsing grammars (IPG) as formalism for expressing some context-sensitive
formats present in binary data, like video and image formats. An IPG attaches to
every non terminal/terminal an interval, which specifies the range of input consumed.
By connecting intervals and attributes, the context-sensitive patterns in file formats
can be well handled. Unlike \emph{typed-peg}, IPGs allow the use semantic actions to
perform arbitrary computations over attributes. While such approach provide extra
power, in our view, it clutters the grammar since the parser specification is not formed only by
grammar rules but also program pieces to handle parsed data.

Another work that proposes a formalism to deal with data formats was developed by Lucks
et. al.~\cite{Lucks17}. Authors devised a new class of languages, named
``calc-regular languages'', which enjoy the property of prefix-freeness that is
present in several binary data formats. In order to provide concise specification
of calc-regular languages, author devised \emph{calc-regular expressions}, an
extension of regular expressions. The parsing of such languages can be made by
finite state machines with accumulators, which are proved equivalent to calc-regular
expressions. Zephyr et. al~\cite{Zephyr21} proposed an extension to PEGs, based on
calc-regular expressions, named calc-PEGs which provide a specific length
operator to PEGs. Authors discuss a $O(n^2)$ algorithm to parse calc-PEGs and
present a tool, called pegmatite, that produce parsers from calc-PEGs specifications.
Unlike our work, the authors do offer no guarantee of termination, neither provide
a detailed specification of calc-PEGs.

The use of intrisincally-typed interpreters in semantics are the subject of
some recent research works~\cite{Feitosa19, Feitosa22, Thiemann23, Rouvoet20,
Rest22}.
Feitosa and Ribeiro~\cite{Feitosa22} present a mechanically verified
implementation of the list-machine benchmark using the dependently typed language
Agda. Authors provide provably correct algorithms for subtyping and least common
supertype calculation, as well as a type checker that produces intrinsically-typed
program representations. They show that this approach reduces code size significantly
— using only 14\% of the lines of code compared to a Twelf solution and 47\% compared
to Coq — while still covering all 14 benchmark tasks. The paper also extends the
formalization to support indirect jumps (version 2.0 of the benchmark), demonstrating
how the intrinsic approach handles more complex control flow features like continuation
types and mutually inductive relations. The formalization of Featherweight Java (FJ)
using intrisincally-typed syntax was the subject of~\cite{Feitosa19}.
Authors used a novel encoding that integrates FJ's nominal type system, class tables,
and inheritance into an intrinsically-typed syntax, where only well-typed expressions
are representable. By using de Bruijn indices for names and splitting class definitions
into signatures and implementations, the authors handle complex binding structures and
mutual dependencies elegantly. They further implement a fuel-based definitional
interpreter that leverages the type information to guarantee progress and preservation
theorems without separate proofs.
Thiemann's work~\cite{Thiemann23} uses a definitional interpreter to define a callback-based approach for
implementing binary session types in dependently typed functional languages like Agda,
eliminating the need for linear types in the host language. By encapsulating all
communication logic within an interpreter that forms the trusted computing base,
the approach guarantees protocol fidelity and linearity by construction. The work extends
to advanced session type features including branching, recursion, multichannel sessions,
higher-order sessions, and context-free sessions, while also introducing novel constructs
such as dynamic selection and an UNROLL command for recursive client protocols.
Another work which uses a definitional interpreter-based semantics for a linear and session-typed
language is~\cite{Rouvoet20}. The authors use a syntax encoding which takes inspiration from separation
logic, the authors develop proof-relevant separation algebras and composable
monadic abstractions (including linear Reader and State monads, and a free
monad for concurrency) that encapsulate the complexities of resource
separation and linearity. These abstractions enable the concise specification of
interpreters for increasingly complex languages: a linear lambda calculus,
one with strong references, and finally a session-typed language with
concurrency and asynchronous communication. The work demonstrates that
with the right abstractions, intrinsically-typed interpreters for sophisticated linear
languages can retain the clarity and simplicity of those for simple languages, while
providing executable and type-safe semantics by construction.

\section{Conclusion}~\label{sec:conclusion}

In this work we proposed a formalization of an extended version of Pest style PEG
grammars. We provide a formalization of stack manipulation and repetition operators
for Pest grammars and presented a mechanized semantics and a soundness result which
implies termination of parsing for all typed grammars.
Besides a formal definition of the operator semantics, we
extended a typing inference algorithm for PEGs~\cite{Cardoso23} which guarantee
termination of well-typed grammars to handle Pest operators.
As future work, we
intend to continue develop of the \emph{typed-peg} Racket library and to
integrate the ideas present in this work in future versions of the Pest tool.

\section*{ACKNOWLEDGMENTS}
This work is supported by the FAPEMIG - Fundação de Amparo à Pesquisa de Minas Gerais -
under grant number APQ-01683-21.


%% bibitems, please use
%%
\bibliographystyle{elsarticle-num}
\bibliography{references}

%% else use the following coding to input the bibitems directly in the
%% TeX file.

%% Refer following link for more details about bibliography and citations.
%% https://en.wikibooks.org/wiki/LaTeX/Bibliography_Management

\end{document}

\endinput
%%
%% End of file `elsarticle-template-num.tex'.
